\section{From Arrays to LAPACK}
\label{sec:lu-packing}

\subsection{The Work Buffer}

stardust의 배열은 임의의 바이트 간격(stride)을 가진다. C 순서(행 우선)가 기본이고, 전치하거나 잘라 낸 뷰는 열 우선이거나 띄엄띄엄하며, 간격이 음수이거나($x[::-1]$) $0$(브로드캐스트된 축)일 수도 있다. LAPACK은 열 우선으로 연속하고 열 사이 간격이 leading dimension $\mathit{ld}$인 행렬만 받는다. 게다가 \code{getrf}는 입력을 덮어쓴다. 그래서 LU 계열의 모든 함수는 행렬을 먼저 작업 버퍼에 옮기며, 이 일을 \term{패킹}{packing}이라 한다. 패킹은 한 번의 패스로 세 가지를 한다. (i) 임의의 간격으로 읽어 열 우선으로 쓴다. (ii) \code{bool}과 \code{int64}를 \code{float64}로 바꾼다. (iii) NaN과 $\pm\infty$가 있는지 본다. LAPACK은 유한하지 않은 입력에 대해 아무것도 약속하지 않으므로(무한 루프나 쓰레기 값이 날 수 있다) stardust는 LAPACK을 부르기 전에 반드시 검사한다.

\begin{algorithm}[t]
\begin{algorithmic}[1]
\Require 원본 시작 $p$, 크기 $m\times n$, 행·열 간격 $r_s$, $c_s$ (바이트), 목적지 $D$와 $\mathit{ld}$
\Ensure $D[j\cdot\mathit{ld}+i]=x_{ij}$ (변환 포함); 모든 원소가 유한한지
\If{$r_s=$ 원소 크기} \Comment{열이 연속: 열마다 복사}
  \For{$j=0,\dots,n-1$}
    \State 열 $j$를 복사하면서 비트로 검사: $\mathit{bad}\gets\mathit{bad}\lor[(w\mathbin{\&}\mathtt{0x7ff0\ldots0})=\mathtt{0x7ff0\ldots0}]$
  \EndFor
\ElsIf{$c_s=$ 원소 크기} \Comment{행이 연속(C 순서): 타일 전치}
  \For{행 타일 $i_0=0,32,64,\dots$} \For{열 타일 $j_0=0,32,64,\dots$}
    \For{$j\in[j_0,j_0+32)$} \For{$i\in[i_0,i_0+32)$}
      \State $D[j\cdot\mathit{ld}+i]\gets x_{ij}$
    \EndFor\EndFor
  \EndFor\EndFor
  \State 써 넣은 열마다 비트 검사를 한 번 더 돈다
\Else \Comment{그 밖: 원소마다}
  \State 열마다, 행마다 읽고 검사하며 쓴다
\EndIf
\State \Return \textbf{not} $\mathit{bad}$
\end{algorithmic}
\caption{패킹 (stardust \code{pack}의 세 경로, 실수 \code{float64}의 경우)}
\label{alg:lu-pack}
\end{algorithm}

\cref{alg:lu-pack}의 경로마다 고른 이유가 있다.
\begin{itemize}
\item \textbf{비트로 하는 유한성 검사.} binary64에서 NaN과 $\pm\infty$는 지수 11비트가 모두 $1$인 수이고 그 밖에는 없다(10장). 그래서 64비트 정수로 읽어 지수 마스크를 비교하고 그 결과를 \code{or}로 모은다. 정수의 \code{or}는 결합적이어서 컴파일러가 루프를 SIMD로 벡터화한다. 부동소수점 비교와 \code{and}로 쓰면 결과는 같지만, stardust가 켠 \code{-ffp-contract=off} 아래에서 컴파일러가 이것을 순서가 정해진 축약으로 보아 벡터화하지 않았다.
\item \textbf{C 순서의 타일 전치.} 행 우선 원본을 열 순서로 그냥 읽으면 원소 하나를 읽을 때마다 행 하나($8n$바이트)를 건너뛴다. $1000\times1000$이면 원소마다 다른 캐시 줄이다. $32\times32$ 타일 하나는 읽는 쪽 32줄과 쓰는 쪽 32줄, 합쳐 $16$KB만 건드리므로 L1 캐시에 머문다(\cref{fig:lu-tile}). 원소가 $4096$개를 넘으면 타일 안을 다시 $4\times4$ 블록으로 나누어 NEON 레지스터 넷에 행 넷을 읽고 \code{trn1}·\code{trn2}로 열 넷을 만든다. 원소 루프의 약 $1.5$배이다. 원소가 $2^{17}$개를 넘는 C 순서 행렬은 열 $32$개 단위로 스레드 풀에 나누어 전치한다. 각 태스크가 자기 열만 쓰므로 결과는 직렬과 같다.
\item \textbf{어느 경로든 같은 결과.} 원소 $(i,j)$는 언제나 $D[j\cdot\mathit{ld}+i]$로 가고 변환은 원소마다 같은 함수이므로, 입력의 배치(C, F, 띄엄띄엄, 음수 간격, 브로드캐스트)가 무엇이든 작업 버퍼는 바이트까지 같다. 그래서 패킹을 거치는 \code{lu}, \code{solve}, \code{inv}, \code{det}, \code{slogdet}의 결과는 입력의 메모리 배치와 무관하게 비트까지 같다.
\end{itemize}

\begin{figure}[t]
\centering
\begin{tikzpicture}[x=0.32cm,y=0.32cm,font=\small]
\draw (0,0) rectangle (12,12); \node at (6,-1) {원본 (행 우선)};
\fill[gray!40] (4,4) rectangle (8,8);
\foreach \r in {4.5,5.5,6.5,7.5}{\draw[->] (4.2,\r) -- (7.8,\r);}
\draw (15,0) rectangle (27,12); \node at (21,-1) {작업 버퍼 (열 우선)};
\fill[gray!40] (19,4) rectangle (23,8);
\foreach \c in {19.5,20.5,21.5,22.5}{\draw[->] (\c,7.8) -- (\c,4.2);}
\draw[->,thick] (8.6,6) -- (18.4,6) node[midway,above] {타일 하나 $32\times32$};
\end{tikzpicture}
\caption{C 순서 원본의 타일 전치. 원본에서는 타일의 행을 따라(화살표) 연속으로 읽고 버퍼에서는 열을 따라 연속으로 쓴다. 한 타일이 건드리는 캐시 줄은 양쪽 합쳐 64줄이다.}
\label{fig:lu-tile}
\end{figure}

패킹은 출력 자리를 작업 버퍼로 쓰기도 한다. \code{solve}는 우변 $b$를 처음부터 결과 배열의 자리에 열 우선으로 패킹하고, 그 자리에서 \code{getrs}(또는 대입)가 $b$를 해로 덮어쓴다. 행렬 $a$만 따로 작업 버퍼가 필요하다. 그래서 \code{solve} 한 번의 메모리 이동은 ``$a$의 복사 한 번, $b$의 복사 한 번, LAPACK''이 전부이다.

\subsection{Leading Dimension and Alignment}

작업 버퍼의 leading dimension은 $\mathit{ld}=n$이다. \cref{sec:lu-blocked}에서 직접 짠 블록 LU가 $\mathit{ld}$가 2의 거듭제곱일 때 느려지는 것을 보았는데 Accelerate의 \code{dgetrf}도 같다. $n=1024$에서 $\mathit{ld}=1024$이면 초당 $149\times10^9$ 플롭, $\mathit{ld}=1032$이면 $235\times10^9$ 플롭이다(\code{figures/48/lda2.c}). stardust는 $\mathit{ld}=n$으로 패킹하므로 $1024\times1024$의 \code{solve}는 이 손해를 그대로 본다. 버퍼의 $\mathit{ld}$를 몇 원소 늘리면 피할 수 있지만 그러면 결과의 끝 비트가 바뀔 수 있어 \cref{sec:lu-testing}의 비트 관계를 함께 다시 확인해야 한다.

정렬은 더 미묘하다. Accelerate는 배열이 16바이트 경계에서 시작하는지에 따라 다른 코드 경로를 타고, 그러면 반올림이 달라진다. \code{float64}의 $3\times3$ 행렬들이 한 배열에 이어 있으면 행렬 사이 간격이 $72$바이트여서 홀수 번째 행렬은 16바이트 경계에서 8바이트 어긋난다(\cref{fig:lu-align}). 같은 행렬이 배치 안의 위치에 따라 다른 비트를 받는 것이다. stardust의 설계 문서는 $n=3,5,7,9$의 무작위 행렬 $120$개에서 \code{schur}, \code{svd}, \code{eigh}의 결과가 수십 개씩 달랐다고 적고 있다. 그래서 LAPACK이나 BLAS에 넘기는 모든 배열을 16바이트 경계에 둔다. 출력 자리가 어긋나 있으면 \code{LapackSlot}이 정렬된 임시 버퍼를 주고, 계산이 끝나면 그 내용을 출력 자리로 복사한다. 그 결과 배치의 어느 자리에 있든 같은 행렬은 같은 비트를 받는다.

\begin{figure}[t]
\centering
\begin{tikzpicture}[x=0.031cm,y=0.5cm,font=\small]
\foreach \k in {0,...,4}{
  \pgfmathsetmacro{\xa}{72*\k}
  \pgfmathtruncatemacro{\odd}{mod(\k,2)}
  \ifnum\odd=1 \fill[gray!35] (\xa,0) rectangle (\xa+72,1); \else \fill[gray!10] (\xa,0) rectangle (\xa+72,1); \fi
  \draw (\xa,0) rectangle (\xa+72,1);
  \node at (\xa+36,0.5) {행렬 \k};
  \node[below] at (\xa,0) {\pgfmathprintnumber{\xa}};
}
\foreach \b in {0,16,...,352}{\draw[gray] (\b,1) -- (\b,1.35);}
\node[anchor=west] at (0,1.75) {16바이트 경계(위의 눈금)};
\end{tikzpicture}
\caption{$3\times3$ \code{float64} 행렬 다섯 개가 이어진 배치. 행렬은 $72$바이트마다 시작하므로 홀수 번째(진한 칸)는 16바이트 경계에서 8바이트 어긋난다.}
\label{fig:lu-align}
\end{figure}

\section{The LU Family in stardust}
\label{sec:lu-stardust-family}

\subsection{Factorization: \texorpdfstring{\code{lu}}{lu}}

\code{sd.linalg.lu(x)}는 $m\times n$ 행렬(과 그 배치)을 받아 \code{LUResult(L, U, perm)}을 돌려준다. $p=\min(m,n)$이면 $L$은 $m\times p$, $U$는 $p\times n$이고 \code{x[..., perm, :] == L @ U}이다. 원소 하나의 몸체는 패킹, \code{getrf}(크기에 따라 \code{getf2}, \cref{sec:lu-batches}), 풀어 쓰기, 순열 만들기의 넷이다. \code{getrf}가 돌려주는 \code{info}$>0$(정확히 $0$인 피벗)은 무시한다. \cref{sec:lu-factorization}에서 보았듯 특이 행렬도 분해는 있고 그 흔적은 $U$의 대각에 남기 때문이다.

풀어 쓰기는 \code{getrf}가 한 자리에 섞어 둔 $L$과 $U$를 두 결과 배열로 나누는 일이다. $L$에는 대각 위에 $0$, 대각에 $1$, 대각 아래에 승수를 쓰고 $U$에는 대각과 그 위를 쓰고 아래에 $0$을 쓴다. 원소가 $1024$개 이하이면 원소마다 조건을 보는 한 루프로 하고, 그보다 크면 열마다 ``$0$ 채우기, $1$, 복사''를 \code{fill}과 \code{copy}로 한다. 작은 행렬에서는 라이브러리 호출 몇 번의 고정비가 원소 루프보다 비싸고, 큰 행렬에서는 연속 구간의 복사가 훨씬 빠르기 때문이다. 마지막으로 \cref{alg:lu-perm}으로 \code{ipiv}를 \code{perm}으로 바꾼다.

\subsection{Solving: \texorpdfstring{\code{solve}}{solve} and \texorpdfstring{\code{lu\_solve}}{lu\_solve}}

\code{solve(a, b)}는 언제나 LU로 푼다. scipy의 \code{assume\_a=}처럼 행렬의 구조를 보고 방법을 고르는 일은 하지 않는다. 대칭 양정치 행렬은 \code{cholesky\_solve}, 삼각 행렬은 \code{solve\_triangular}처럼 구조는 함수 이름으로 고르는 것이 stardust의 규칙이다. $b$가 1차원이면 벡터 하나이고 그 밖에는 $(\dots,n,k)$의 행렬 묶음이다.

\begin{algorithm}[t]
\begin{algorithmic}[1]
\Require 배치 원소 하나의 $a$ ($n\times n$), $b$ ($n\times k$)
\Ensure 결과 자리에 $x=a^{-1}b$, 또는 실패 코드
\State $a$를 작업 버퍼에 패킹; 유한하지 않으면 \Return \code{NonFinite}($a$)
\State $\mathit{info}\gets$ \code{getrf} 또는 \code{getf2}$(a)$ \Comment{$32\le n\le46$이면 \code{getf2}}
\State \textbf{if} $\mathit{info}>0$ \textbf{then return} \code{Singular}
\State $b$를 결과 자리에 패킹; 유한하지 않으면 \Return \code{NonFinite}($b$)
\If{$k\le2$}
  \State 결과의 열마다 \code{ipiv}의 교환을 적용
  \State \code{trsv}(단위 하삼각 $L$), \code{trsv}($U$)를 열마다
\Else
  \State \code{getrs}('N') \Comment{교환 + \code{trsm} 둘}
\EndIf
\end{algorithmic}
\caption{\code{solve}의 배치 원소 하나}
\label{alg:lu-solve-body}
\end{algorithm}

\cref{alg:lu-solve-body}의 순서는 에러의 순서이기도 하다. $a$가 NaN을 품고 있으면 $b$에 무엇이 있든 $a$에 대한 \code{ValueError}이고, $a$가 특이하면 $b$를 보기 전에 \code{LinAlgError}이다. 6--7행이 \code{getrs} 대신 \code{trsv}를 부르는 이유는 측정이다. Accelerate의 \code{trsm}은 우변이 하나일 때 \code{trsv}보다 크게 느리다. \cref{tab:lu-trsv}처럼 $n=1600$에서 세 배이다. \code{getrs}는 안에서 \code{trsm}을 부르므로 우변이 적을 때는 교환과 \code{trsv} 두 번을 직접 하는 편이 빠르다. 우변이 넷 이상이면 \code{trsm}이 블록으로 재사용을 얻어 앞서고, 경계는 $k\le2$에 두었다.

\begin{table}[t]
\begin{sdtable}{colspec={rrrrrrr}, row{1}={font=\small\bfseries}}
\toprule
$n$ & $50$ & $100$ & $200$ & $400$ & $800$ & $1600$ \\
\midrule
\code{trsv} & $0.42$ & $1.50$ & $5.04$ & $8.54$ & $37.4$ & $112$ \\
\code{trsm} & $0.54$ & $3.33$ & $13.0$ & $22.3$ & $100$ & $342$ \\
\bottomrule
\end{sdtable}
\caption{우변 하나의 하삼각 풀이(Accelerate, 마이크로초, \code{figures/48/bench.c}의 \code{trsv}).}
\label{tab:lu-trsv}
\end{table}

\code{lu\_solve(l, u, perm, b, *, trans='N')}는 \code{lu}의 결과를 다시 받아 푼다. scipy의 \code{lu\_factor}가 돌려주는 \code{(lu, piv)} 같은 압축 형식이 아니라 \code{lu}의 결과 그대로를 받는다는 점이 다르고, 그래서 \code{perm}이 정말 순열인지부터 확인해야 한다.

\begin{algorithm}[t]
\begin{algorithmic}[1]
\Require $l$ ($n\times n$, 대각 아래만 읽음), $u$ ($n\times n$, 대각과 위만 읽음), $\mathit{perm}$, $b$ ($n\times k$), $\mathit{trans}\in\{\mathrm N,\mathrm T,\mathrm H\}$
\State $l$의 대각 아래가 유한한지; 아니면 \Return \code{NonFinite}($l$)
\State $u$의 대각과 위가 유한한지; 아니면 \Return \code{NonFinite}($u$)
\State $u$의 대각에 정확한 $0$이 있으면 \Return \code{Singular}
\State $\mathit{seen}_{0:n}\gets0$
\For{$r=0,\dots,n-1$} \Comment{순열 검증}
  \State $v\gets\mathit{perm}_r$; \textbf{if} $v<0$ \textbf{or} $v\ge n$ \textbf{or} $\mathit{seen}_v$ \textbf{then return} \code{BadPerm}
  \State $\mathit{seen}_v\gets1$
\EndFor
\If{$\mathit{trans}=\mathrm N$}
  \State $x_{r,:}\gets b_{\mathit{perm}_r,:}$ (유한성 검사 포함) \Comment{행 모으기: 반올림 없음}
  \State $x\gets L^{-1}x$ (단위 대각); $x\gets U^{-1}x$
\Else
  \State $z\gets b$ (유한성 검사 포함)
  \State $z\gets\mathrm{op}(U)^{-1}z$; $z\gets\mathrm{op}(L)^{-1}z$ \Comment{$\mathrm{op}$은 전치 또는 켤레 전치}
  \State $x_{\mathit{perm}_r,:}\gets z_{r,:}$ \Comment{흩기}
\EndIf
\end{algorithmic}
\caption{\code{lu\_solve}의 배치 원소 하나}
\label{alg:lu-lusolve}
\end{algorithm}

\cref{alg:lu-lusolve}의 전치 경우가 맞는 까닭은 이렇다. \code{perm}의 뜻에서 $PA=LU$이고 $(PA)_r=A_{\mathit{perm}_r}$이다. 그러면 $A=P^{\mathsf T}LU$이고 $A^{\mathsf T}=U^{\mathsf T}L^{\mathsf T}P$이다. $A^{\mathsf T}x=b$는 $U^{\mathsf T}w=b$, $L^{\mathsf T}z=w$, $Px=z$로 풀리고 마지막 식은 $x_{\mathit{perm}_r}=z_r$, 곧 흩기이다. \code{'N'}이면 반대로 $Pb$를 먼저 만드는데 이것은 $b$의 행을 \code{perm}의 순서로 모으는 것이어서 산술이 없다. 그 뒤의 두 삼각 풀이는 \code{solve}가 교환 뒤에 하는 것과 같은 호출이다. 그래서 $l$, $u$가 \code{lu}가 돌려준 열 우선 배열 그대로이면 \code{lu\_solve(*lu(a), b)}는 \code{solve(a, b)}와 비트까지 같다.

$l$과 $u$는 LAPACK이 읽기만 하므로 복사하지 않고 그 자리에서 읽을 수 있다. 원소 타입이 계산 타입과 같고, 시작이 16바이트 경계이고, 열 우선이면서 열 간격이 16바이트의 배수이면 그대로 \code{trsv}·\code{trtrs}에 넘기고 읽을 삼각형만 유한성을 검사한다. C 순서(행 우선)인 삼각 행렬도 그 자리에서 읽는다. 행 우선의 하삼각은 열 우선으로 본 상삼각의 전치이므로, 우변이 둘 이하이면 \code{trsv}에 \code{CblasRowMajor}를 주고 셋 이상이면 \code{trtrs}에 반대쪽 삼각과 반대 연산(\code{'N'}$\leftrightarrow$\code{'T'})을 준다. 다만 이 길은 반올림 순서가 달라 열 우선 인수와 끝 비트가 다를 수 있다. 복소수의 켤레 전치(\code{'H'})에 우변이 셋 이상인 경우는 대응하는 연산이 없어서 패킹한다.

\subsection{Inverse and Determinants}

\code{inv(x)}는 \code{getrf} 뒤에 단위행렬을 우변으로 두고 \cref{alg:lu-solve-body}의 풀이를 그대로 한다. \code{getri}를 쓰지 않는 이유는 둘이다. 하나는 \cref{tab:lu-inverse}에서 보았듯 $n\le1024$에서 이 길이 빠르거나 같다는 것이고, 다른 하나는 \code{inv(a)}가 \code{solve(a, I)}와 비트까지 같아진다는 것이다. 두 함수가 같은 행렬에 대해 반올림 수준에서 다른 답을 주면 사용자가 의아해한다. 대신 $n=2048$처럼 큰 행렬에서는 \code{getri}보다 느리다. 설계 문서는 크기에 따라 길을 바꾸는 안을 남겨 두었다.

\code{det}과 \code{slogdet}는 \code{getrf} 뒤에 \cref{alg:lu-det}과 로그 합을 한다. 특이(\code{info}$>0$)도 에러가 아니다. 빈 행렬($0\times0$)의 행렬식은 빈 곱의 관례대로 $1$, \code{slogdet}은 $(1,0)$이다. 배치가 없으면 결과는 Python 스칼라이고 배치가 있으면(크기가 1이어도) 배열이다.

\subsection{Errors}

LU 계열이 내는 에러는 둘이다. 입력이 NaN이나 무한대를 품으면 \code{ValueError}이고 특이 행렬로 풀거나 역행렬을 구하면 \code{LinAlgError}이다. 메시지는 함수 이름으로 시작한다.
\begin{itemize}
\item \code{solve(): 'a' contains NaN or infinity}
\item \code{inv(): singular matrix (at batch index (1, 2))}
\end{itemize}
배치의 여러 원소가 실패하면 가장 작은 배치 인덱스를 보고한다. 한 원소 안에서는 인자의 순서대로 본다. 첫 행렬의 유한성, 그 행렬의 특이성, 다음 인자들의 유한성이다(\cref{alg:lu-solve-body,alg:lu-lusolve}). ``특이''는 정확히 $0$인 피벗만을 뜻한다. 피벗이 $10^{-300}$이어도 분해는 성공하고 답은 거대하거나 무한대가 되며, 조건수를 보는 것은 사용자의 몫이다(\code{cond}). 이 규칙은 numpy와 같다.

\section{Batches, Threads, and Hardware}
\label{sec:lu-batches}

\subsection{The Batch Driver}

모든 행렬 함수는 앞쪽 축을 \term{배치}{batch}로 본다. \code{solve(a, b)}에서 $a$가 $(2,1,n,n)$이고 $b$가 $(3,n,k)$이면 배치 축 $(2,1)$과 $(3,)$을 오른쪽 정렬로 브로드캐스트해 $(2,3)$개의 문제를 푼다. 원소마다 각 피연산자의 바이트 오프셋을 계산하고(브로드캐스트되는 축의 간격은 $0$) 같은 몸체를 부른다. 배치 원소는 서로 독립이므로 스레드에 나누어 돌릴 수 있다.

\begin{algorithm}[t]
\begin{algorithmic}[1]
\Require 배치 크기 $N$, 원소 하나의 추정 시간 $t$ (ns), 몸체 $\mathit{body}(i)$ (실패 코드 또는 $0$), 태스크 상한 $c$
\If{$N\ge2$ \textbf{and} $Nt\ge10^4$ \textbf{and} 풀이 쓸 수 있음}
  \State $q\gets\max(1,\lfloor5000/t\rfloor)$ \Comment{태스크 하나가 약 $5\,\mu$s}
  \State \textbf{if} $c>0$ \textbf{and} $\lceil N/q\rceil>c$ \textbf{then} $q\gets\lceil N/c\rceil$
  \State $\mathit{first}\gets\infty$ (원자적)
  \ForAll{태스크 $\tau=0,\dots,\lceil N/q\rceil-1$ \textbf{in parallel}}
    \State 자기 작업 메모리를 만든다
    \For{$i=\tau q,\dots,\min(N,(\tau+1)q)-1$}
      \State \textbf{if} $i>\mathit{first}$ \textbf{then break} \Comment{더 작은 실패가 이미 있다}
      \If{$\mathit{body}(i)\ne0$}
        \State 기록 $(i,\text{코드})$; $\mathit{first}\gets\min(\mathit{first},i)$ (비교 후 교환); \textbf{break}
      \EndIf
    \EndFor
  \EndFor
  \State \Return 기록들 가운데 인덱스가 가장 작은 것
\Else
  \State GIL을 풀고(오래 걸리면) $i=0,1,\dots$ 차례로 $\mathit{body}(i)$; 첫 실패에서 멈추고 \Return
\EndIf
\end{algorithmic}
\caption{배치 드라이버 (\code{run\_batch})}
\label{alg:lu-batch}
\end{algorithm}

\cref{alg:lu-batch}에서 원소 하나의 추정 시간은 $t=a+c\,w$의 꼴이다. $w$는 행렬의 원소 수이고 LU 계열은 $a=70$ns, $c=3$ns로 $n=2..256$의 실측에 맞춘 상수이다. $n$이 크면 시간이 $n^3$으로 자라 이 식은 과소평가하지만, 그때는 원소 둘만 있어도 이미 나눌 가치가 있으므로 판단은 틀리지 않는다. 배치 전체의 추정 시간이 $10\,\mu$s 미만이면 스레드를 깨우는 비용이 이득보다 크므로 나누지 않는다.

실패 보고는 직렬과 같아야 한다. 직렬로 돌면 처음 만난 실패가 가장 작은 인덱스이다. 나누어 돌릴 때 처음 발견한 실패에서 모두 멈추면 어느 인덱스가 보고될지가 스레드의 타이밍에 달린다. \cref{alg:lu-batch}은 이렇게 한다. 각 태스크는 자기 구간을 앞에서부터 돌다 첫 실패에서 멈추고 그 인덱스를 기록한다. 지금까지 알려진 가장 작은 실패보다 뒤의 원소는 건너뛴다. 끝나면 기록들 가운데 가장 작은 것을 고른다. 가장 작은 실패 인덱스 $i^\ast$를 품은 태스크는 $i^\ast$에 이르기 전에 멈추지 않는다. 그 앞의 원소들은 실패하지 않으므로 기록되지 않고, $\mathit{first}$는 언제나 어떤 실제 실패의 인덱스여서 $i^\ast$ 이상이므로 건너뛰기 조건도 걸리지 않는다. 그러므로 $i^\ast$는 반드시 기록되고, 보고되는 인덱스는 스레드 수와 상관없이 직렬과 같다.

결과의 비트도 스레드 수와 상관없다. 원소 하나는 어디서 돌든 같은 몸체를 같은 순서로 돌고, 작업 메모리는 태스크마다 따로이며, LAPACK에 넘기는 배열은 모두 16바이트 경계에 있기 때문이다. 작업 메모리에도 두 가지 세부가 있다. 하나는 새 메모리를 $0$으로 채우지 않는 것이다. \code{std::vector}의 \code{resize}처럼 채우면 $1000\times1000$에서 그 비용이 패킹만큼 들었고, 모든 사용처가 읽기 전에 쓰므로 채울 필요가 없다. 다른 하나는 $4$KB보다 작은 버퍼를 $128$바이트 경계(이 처리기의 캐시 줄)에 두고 크기도 $128$의 배수로 올리는 것이다. 그러지 않으면 메모리 할당기가 여러 스레드의 작은 버퍼(\code{lu\_solve}의 $3$바이트짜리 \code{seen})를 같은 캐시 줄에 내주는 일이 생겼고, 원소마다 쓰는 그 줄이 코어 사이를 오가며 시간이 프로세스마다 30\%씩 달라졌다.

\subsection{A Shared Matrix Unit}

Apple의 처리기에는 행렬 곱을 위한 장치(AMX)가 코어 무리마다 하나씩 있고 Accelerate의 일부 경로가 이것을 쓴다. 이 장치는 코어들이 함께 쓰므로 여러 스레드가 동시에 그 경로를 타면 줄을 선다. \cref{fig:lu-amx}는 크기가 같은 행렬을 수천 개 분해하는 시간을 행렬 하나당으로 나타낸 것이다. 한 스레드에서는 \code{getrf}와 \code{getf2}가 비슷하다. 열 스레드에서 \code{getf2}는 $n$을 따라 고르게 늘지만 \code{getrf}는 $n=32$에서 갑자기 $1.5\,\mu$s로 뛰고 $n=64$까지 그 자리에 머문다. 스레드 수를 늘려도 빨라지지 않는 이 평평한 구간이 공유 장치에서 줄을 서는 모습이다. $n=36$에서 \code{getf2}는 \code{getrf}의 $4.2$배 빠르다.

\begin{figure}[t]
\centering
\begin{tikzpicture}
\begin{axis}[width=0.86\textwidth,height=0.5\textwidth,
  xlabel={$n$},ylabel={행렬 하나당 시간 ($\mu$s)},grid=major,grid style={gray!25},
  legend pos=north west,legend style={font=\small,draw=none,fill=none},
  xmin=6,xmax=74,ymin=0,ymax=2.5]
\addplot[black,mark=*,mark size=1.2pt] table[x=n,y=getrf10]{figures/48/batch.dat};
\addlegendentry{\code{getrf}, 10스레드}
\addplot[black,dashed,mark=square*,mark size=1.2pt] table[x=n,y=getf210]{figures/48/batch.dat};
\addlegendentry{\code{getf2}, 10스레드}
\draw[gray!60] (axis cs:32,0) -- (axis cs:32,2.5);
\draw[gray!60] (axis cs:46,0) -- (axis cs:46,2.5);
\node[font=\scriptsize,anchor=south] at (axis cs:39,2.05) {stardust의 \code{getf2} 구간};
\end{axis}
\end{tikzpicture}
\caption{$n\times n$ 행렬 수천 개를 열 스레드로 분해할 때 행렬 하나당 시간(Apple M4 Pro, Accelerate). 한 스레드의 값과 측정 코드는 \code{figures/48/batch.dat}과 \code{bench.c}에 있다.}
\label{fig:lu-amx}
\end{figure}

stardust는 이 측정을 두 규칙으로 바꿨다. 첫째, 행과 열이 모두 $32$에서 $46$ 사이인 행렬은 \code{getrf} 대신 블록 없는 \code{getf2}로 분해한다. 피벗 규칙과 특이의 판정은 같고 결과는 반올림 수준에서 같다. 둘째, 우변이 $8$개 이상인 삼각 풀이(\code{trsm})도 $n=24$--$40$에서 같은 장치를 타므로 이 경우 배치를 나누는 태스크의 수를 $3$이나 $5$로 제한한다(\cref{alg:lu-batch}의 $c$). 경계값은 stardust를 개발한 기계에서 잰 것이다. 우리 측정에서는 열 스레드의 \code{getf2}가 $n=64$까지 앞서므로 경계는 처리기마다 다르다. 이런 규칙은 하드웨어를 바꿀 때마다 다시 재야 한다는 점에서 알고리즘의 일부가 아니라 구현의 일부이다.

\section{Testing and Accuracy}
\label{sec:lu-testing}

후진 오차 해석은 시험의 기준이 된다. stardust는 풀이의 \term{잔차 계약}{residual contract}을
\[
\|Ax-b\|_\infty\le c\,n\,\epsilon\,(\|A\|_\infty\|x\|_\infty+\|b\|_\infty),\qquad c=10
\]
으로 둔다. \cref{thm:lu-rigal-gaches}의 후진 오차를 $\infty$-노름으로 쓴 것에 $c\,n\,\epsilon$의 상한을 준 꼴이다. \eqref{eq:lu-normwise}의 이론 상한 $n^2\rho_n\gamma_{3n}$은 너무 비관적이어서 시험 기준으로 쓸 수 없고, 상수 $c=10$은 측정에서 정했다. \cref{tab:lu-residual}은 무작위 정규 행렬에서 잰 이 비율이다. 가장 나쁜 값도 계약의 $1/90$ 아래이다.

\begin{table}[t]
\begin{sdtable}{colspec={rrrr}, row{1}={font=\small\bfseries}}
\toprule
$n$ & 행렬 수 & 중앙값 & 최댓값 \\
\midrule
$10$ & $2000$ & $0.019$ & $0.112$ \\
$100$ & $300$ & $0.0055$ & $0.0092$ \\
$1000$ & $20$ & $0.0036$ & $0.0049$ \\
\bottomrule
\end{sdtable}
\caption{$\|Ax-b\|_\infty/(n\epsilon(\|A\|_\infty\|x\|_\infty+\|b\|_\infty))$, 표준 정규 원소의 $A$와 $b$, 부분 피벗 LU(Accelerate)로 푼 해(\code{figures/48/residual.py}). 계약은 이 값이 $10$ 이하일 것을 요구한다.}
\label{tab:lu-residual}
\end{table}

이 밖에 stardust의 시험은 다음을 확인한다.
\begin{itemize}
\item \textbf{조건수가 큰 행렬}: $\kappa_2=1$에서 $10^{12}$까지의 정사각 행렬에 대해 해와 행렬식의 참값을 구간 산술 라이브러리 Arb로 미리 계산해 두고, 잔차 계약과 함께 확인한다.
\item \textbf{비트 관계}: \code{inv(a) == solve(a, I)}, \code{lu\_solve(*lu(a), b) == solve(a, b)}(열 우선 인수), \code{matrix\_power(x, -1) == inv(x)}가 바이트까지 같은지.
\item \textbf{배치와 레이아웃}: 같은 행렬을 C, F, 띄엄띄엄, 음수 간격으로 넣었을 때, 배치의 홀수 자리($n=3,5,7,40$)에 두었을 때, 스레드를 $1$, $3$, $10$개로 바꿨을 때 결과가 같은 바이트인지.
\item \textbf{경계 크기}: \code{getf2} 구간의 양 끝 $n=31,32,46,47$에서 답, 비트 재현성, 특이의 판정.
\item \textbf{규칙}: 빈 행렬, 배치 크기 $0$, NaN 입력의 \code{ValueError}, 특이 행렬의 \code{LinAlgError}와 가장 작은 배치 인덱스, 에러의 순서, 넘치는 부분 곱이 있는 \code{det}(\code{diag(1e200, 1e200, 1e-300)}이 $10^{100}$), 복소 \code{lu}의 $|\ell_{ij}|\le\sqrt2$.
\end{itemize}

numpy와의 비교는 허용오차로 한다. Accelerate의 \code{gesv}는 \code{getrf}와 \code{getrs}를 따로 부를 때와 내부 계산 순서가 달라 $n\ge4$에서 끝 비트가 다르고, numpy의 \code{solve}는 \code{gesv}를 부르기 때문이다. 같은 기계와 같은 스레드 수에서 실행할 때마다 같은 비트를 내는 것이 stardust의 약속이고, 다른 라이브러리와 비트가 같은 것은 약속이 아니다.

\section{APIs in This Chapter}
\label{sec:lu-apis}

\begin{table}[h]
\begin{sdtable}{colspec={lXl}, row{1}={font=\small\bfseries}}
\toprule
API & 하는 일 & 설명한 곳 \\
\midrule
\code{lu(x)} & $x[\mathit{perm}]=LU$, 직사각·특이 허용 & \cref{sec:lu-pivoting,sec:lu-blocked,sec:lu-stardust-family}, \cref{alg:lu-blocked,alg:lu-recursive,alg:lu-perm} \\
\code{solve(a, b)} & $ax=b$, 배치·브로드캐스트 & \cref{sec:lu-using,sec:lu-stardust-family}, \cref{alg:lu-solve,alg:lu-solve-body} \\
\code{lu\_solve(l, u, perm, b, *, trans)} & \code{lu}의 결과로 풀기, 전치 & \cref{sec:lu-stardust-family}, \cref{alg:lu-lusolve} \\
\code{inv(x)} & 역행렬, \code{solve(x, I)}와 같은 비트 & \cref{sec:lu-using,sec:lu-stardust-family}, \cref{tab:lu-inverse} \\
\code{det(x)} & 넘침을 피하는 행렬식 & \cref{sec:lu-using}, \cref{alg:lu-det} \\
\code{slogdet(x)} & 부호와 $\log|\det|$ & \cref{sec:lu-using} \\
\bottomrule
\end{sdtable}
\caption{이 장이 구현까지 설명한 \code{sd.linalg}의 API. 패킹과 배치 드라이버(\cref{sec:lu-packing,sec:lu-batches})는 이 장의 모든 API가 함께 쓴다.}
\label{tab:lu-apis}
\end{table}

\section{Concluding Remarks}

가우스 소거를 행렬의 분해로 보는 순간 알고리즘은 세 층으로 갈린다. 산술의 층에서는 승수를 제자리에 적는 것만으로 $L$이 생기고, 부분 피벗과 성장 인자가 안정성을 정한다. 이론 상한 $2^{n-1}$과 실제 성장 $n^{2/3}$ 사이의 간극은 아직 완전히 설명되지 않았다. 메모리의 층에서는 같은 산술을 \code{gemm}으로 바꾸어 데이터 하나를 여러 번 쓰게 하는 것이 속도를 열 배 넘게 바꾸며, 블록과 재귀가 그 방법이다. 기계의 층에서는 leading dimension의 모양, 16바이트 정렬, 코어들이 함께 쓰는 행렬 장치처럼 알고리즘 책에 나오지 않는 것들이 속도와 비트 재현성을 좌우한다. stardust의 LU 계열은 산술을 LAPACK에 맡기는 대신 나머지 두 층에서 패킹, 정렬, 배치 드라이버, 크기별 경로라는 형태로 일을 한다. 다음 장에서는 삼각 행렬과 띠 행렬처럼 구조가 있는 연립방정식을 다룬다.

LU와 그 오차 해석의 표준 참고서는 Higham의 책 9장과 Golub과 Van Loan의 책 3장이다\cite{higham2002accuracy,golub2013matrix}. 가우스 소거의 역사는 Grcar가 자세히 정리했다\cite{grcar2011}.

\chapterreferences
